% 编译入口：xelatex main.tex（至少两遍），或 tectonic main.tex。
% 这是个人结构与版式模板，不是已完成或可直接提交的论文。
\PassOptionsToPackage{fontset=none}{ctex}
\documentclass[withoutpreface]{cumcmthesis}
% 从现有最终论文提取的个人版式；与题目内容无关的样式集中在此。
% 原正式版本中的按页 Needspace、图片裁剪等局部补丁未复制。
% 字体优先保持原稿在macOS上的宋体/黑体；无此字体时使用TeX发行版Fandol。
\IfFontExistsTF{Songti SC}{
  \setCJKmainfont{Songti SC}[BoldFont=Songti SC Bold]
  \setCJKsansfont{Heiti SC}
  \setCJKmonofont{Songti SC}
  \setCJKfamilyfont{zhsong}{Songti SC}
  \setCJKfamilyfont{zhhei}{Heiti SC}
  \setCJKfamilyfont{zhkai}{Songti SC}
  \setCJKfamilyfont{zhfs}{Songti SC}
}{
  \setCJKmainfont{FandolSong-Regular.otf}[BoldFont=FandolSong-Bold.otf]
  \setCJKsansfont{FandolHei-Regular.otf}[BoldFont=FandolHei-Bold.otf]
  \setCJKmonofont{FandolSong-Regular.otf}
  \setCJKfamilyfont{zhsong}{FandolSong-Regular.otf}
  \setCJKfamilyfont{zhhei}{FandolHei-Regular.otf}
  \setCJKfamilyfont{zhkai}{FandolKai-Regular.otf}
  \setCJKfamilyfont{zhfs}{FandolFang-Regular.otf}
}
\providecommand{\songti}{\CJKfamily{zhsong}}
\providecommand{\heiti}{\CJKfamily{zhhei}}
\providecommand{\kaishu}{\CJKfamily{zhkai}}
\providecommand{\fangsong}{\CJKfamily{zhfs}}
\usepackage{needspace}
\usepackage{fvextra}
\usepackage{placeins}
\usepackage{microtype}
\DeclareMathSizes{9.5}{9.5}{8}{8}
\newcommand{\important}[1]{{\heiti\bfseries #1}}
\newcommand{\questionlabel}[1]{{\heiti\bfseries\fontsize{14}{18}\selectfont #1}}
\newcommand{\topic}[1]{{\heiti\bfseries\fontsize{11}{14}\selectfont #1}}
\newcommand{\questionline}[1]{\par\noindent\hspace*{2em}{\heiti\bfseries\fontsize{14}{18}\selectfont #1\par}}
\setmainfont{texgyretermes-regular.otf}[BoldFont=texgyretermes-bold.otf,ItalicFont=texgyretermes-italic.otf,BoldItalicFont=texgyretermes-bolditalic.otf]
\setsansfont{texgyreheros-regular.otf}[BoldFont=texgyreheros-bold.otf,ItalicFont=texgyreheros-italic.otf,BoldItalicFont=texgyreheros-bolditalic.otf]
\IfFontExistsTF{Menlo}{\setmonofont{Menlo}}{\setmonofont{lmmono10-regular.otf}}
\xeCJKDeclareCharClass{CJK}{"22A5}
\geometry{a4paper,margin=25mm}
\renewcommand{\baselinestretch}{1.28}
\renewcommand{\thesection}{\chinese{section}、}
\renewcommand{\thesubsection}{\arabic{section}.\arabic{subsection}}
\renewcommand{\thefigure}{1-\arabic{figure}}
\renewcommand{\thetable}{1-\arabic{table}}
\makeatletter
\renewcommand{\mcm@cap@keywordsname}{关键词}
\renewcommand\keywords[1]{%
  \renewcommand{\mcm@tokens@keywords}{#1}%
  \par\vskip1ex
  {\noindent\zihao{-4}\heiti\bfseries\mcm@cap@keywordsname：}~{\mcm@tokens@keywords}%
}
\renewcommand\section{\@startsection{section}{1}{\z@}%
 {-3.0ex \@plus -1ex \@minus -.2ex}{1.8ex \@plus .2ex}%
 {\centering\normalfont\heiti\fontsize{14}{18}\selectfont\bfseries}}
\renewcommand\subsection{\@startsection{subsection}{2}{\z@}%
 {-1.4ex \@plus -.5ex \@minus -.2ex}{.5ex \@plus .1ex}%
 {\normalfont\heiti\fontsize{14}{18}\selectfont\bfseries}}
\renewcommand\subsubsection{\@startsection{subsubsection}{3}{2em}%
 {-1.5ex \@plus -.5ex \@minus -.2ex}{.7ex \@plus .2ex}%
 {\normalfont\heiti\fontsize{11}{14}\selectfont\bfseries}}
\makeatother
\captionsetup[figure]{font={small,normalfont},labelfont={bf},textfont={normalfont},skip=5pt}
\captionsetup[table]{font={small,normalfont},labelfont={bf},textfont={normalfont},skip=5pt}
\widowpenalty=10000 \clubpenalty=10000 \displaywidowpenalty=10000
\makeatletter
\setlength{\@fptop}{0pt}\setlength{\@fpsep}{14pt}\setlength{\@fpbot}{0pt plus 1fil}
\makeatother
\setlength{\textfloatsep}{12pt plus 3pt minus 2pt}
\setlength{\intextsep}{12pt plus 3pt minus 2pt}
\setlength{\emergencystretch}{2em}
\pagestyle{plain}
% 每问图表采用“问号-序号”；总路线图延续原稿计入第一问图序。
\newcommand{\startquestion}[3]{%
  \FloatBarrier
  \subsection{问题#2}
  \renewcommand{\thefigure}{#1-\arabic{figure}}%
  \renewcommand{\thetable}{#1-\arabic{table}}%
  \ifnum#1>1\setcounter{figure}{0}\fi
  \setcounter{table}{0}%
  \questionline{第#2问：#3}%
}
% 仅用于模板说明，不自动进入未来正式论文。
\newcommand{\templatehint}[1]{\noindent{\small\textcolor{gray}{［模板说明］#1}}\par}

% 保留短题名单行下划线；长题名自然换行，避免不可断行下划线盒溢出。
\newlength{\personalTitleWidth}
\newcommand{\papertitlecontent}{%
  \settowidth{\personalTitleWidth}{\PaperTitle}%
  \ifdim\personalTitleWidth>\linewidth
    \PaperTitle\par\nobreak\vskip3pt\hrule
  \else
    \underline{\PaperTitle}%
  \fi
}

% 集中修改身份与题名。正式成稿时删除“模板样张”等说明。
\newcommand{\PaperTitle}{数学建模论文题目（模板样张）}
\newcommand{\TeamCode}{待填队伍编号}
\newcommand{\TrackName}{待填赛道或题号}
\newcommand{\PaperKeywords}{关键词一；关键词二；关键词三；关键词四}
% true 显示其格式位置；正式使用时按当届要求决定，不继承旧比赛AI历史。
\newif\ifshowdisclosure \showdisclosuretrue
\newif\ifshowappendix \showappendixtrue
% 题数由 sections/questions.tex 的输入清单决定；默认为四问格式示例。

\hypersetup{pdftitle={\PaperTitle},pdfauthor={},pdfsubject={个人数学建模论文结构与格式模板},pdfkeywords={\PaperKeywords},hidelinks,pageanchor=true}
\title{\PaperTitle}
\date{}
\begin{document}
\pagenumbering{roman}
\pagestyle{empty}
\noindent{\heiti\fontsize{10}{12}\selectfont 队伍编号：\TeamCode\hfill 赛道／题号：\TrackName}\par
\begin{center}
{\songti\bfseries\fontsize{16}{21}\selectfont
\begin{tabularx}{\textwidth}{@{}l@{\quad}X@{}}
题\quad 目 & \papertitlecontent
\end{tabularx}}
\end{center}
\begin{center}
{\heiti\bfseries\fontsize{14}{18}\selectfont 摘\quad 要：}
\end{center}
本文档为当前个人论文结构与版式的可编译样张。这里填写比赛背景、研究对象与总体研究主线，用一段话连接各问，避免把题面原文完整搬入摘要。

\questionlabel{针对问题一，}概括本问需要解决的问题、实际采用的方法及最重要的结果。模型名称可用\important{黑体加粗}突出；正式数值应从已确认结果表转录，模板不提供研究结果。

\questionlabel{针对问题二，}说明本问如何承接前问的有效信息，交代模型的主要数学关系和关键结论。只写实际完成的步骤，不把前问排序直接变成物理参数。

\questionlabel{针对问题三，}概括方法设计、求解与验证方式，给出有对应证据的答案。涉及区间、概率或情景时，应与正文采用相同定义。

\questionlabel{针对问题四，}总结本问的主要贡献及关键结果，并与整篇研究主线衔接。若实际题数不同，摘要和正文的问题输入清单需一起调整。

\templatehint{本页仅演示当前摘要结构、字体与重点标记。正式稿删除所有模板说明，替换为本次比赛内容。}

\keywords{\important{\PaperKeywords}}
\clearpage
\setcounter{tocdepth}{2}
\setlength{\cftsecnumwidth}{3em}
\tableofcontents
\thispagestyle{empty}
\clearpage
\pagenumbering{arabic}
\pagestyle{plain}
\section{问题重述}
\subsection{问题背景}
本节填写理解赛题所必需的背景、研究对象及实际问题。先交代对象和已有条件，再引出研究任务；引用真实来源时使用正文引用标记。示例文献\cite{template-guide}仅用于展示参考文献编号，不是研究依据。

\subsection{问题的提出}
（1）填写第一问的输入、已知条件和需要给出的答案。

（2）填写第二问相对第一问新增的条件、约束或研究对象。

（3）填写第三问的任务、评价要求及需要回答的核心问题。

（4）填写第四问的目标及与前面各问的关系。

\templatehint{四问是现有合稿的结构示例。后续按实际题目增减，不把题数或问间依赖固定为比赛要求。}

\section{问题分析}
\subsection{问题分析}
先从整题角度说明各问怎样构成一条研究主线，区分共用资料与每问独有输入；下文再逐问解释本问困难、数据条件及方法思路。

\questionlabel{针对问题一，}填写研究目标、可直接使用的资料、需要构造的变量，以及为什么选择相应方法。

\questionlabel{针对问题二，}填写对前问信息的实际依赖，并说明新增条件怎样改变模型或求解目标。

\questionlabel{针对问题三，}填写需要进一步解决的疑点、候选路线和相应验证依据。

\questionlabel{针对问题四，}填写本问的建模与检验思路，以及如何综合前问有效成果。

需要用方法图解释整题关系时，可替换图~\ref{fig:roadmap}的占位框；无此需要时删除整个图环境，并同步正文引用。
\begin{figure}[htbp]
\centering
\fbox{\parbox[c][18mm][c]{0.88\textwidth}{\centering 图形位置示例：研究对象、方法与输出之间的关系\\此框不是计算结果或正式技术路线图}}
\caption{整题研究技术路线（格式位置示例）}\label{fig:roadmap}
\end{figure}
\FloatBarrier

\section{模型假设}
\begin{enumerate}
\item 填写真正影响答案的第一条假设，并说明采用理由及适用对象。
\item 填写必要的数据支持、边界或信息可得性假设，不将未知量改写成已知事实。
\item 对重要歧义比较合理解释，明确当前采用哪一种及其对后问的影响。
\end{enumerate}
假设的具体应用在相应模型处说明；模型评价部分集中讨论它们造成的适用限制。实际没有采用的假设不因模板中有条目而保留。

\section{符号说明}
\renewcommand{\thetable}{\arabic{table}}
\begin{table}[htbp]
\centering
\caption{主要符号及含义（格式示例）}\label{tab:symbols}
\fontsize{9.5}{12}\selectfont
\setlength{\tabcolsep}{4pt}
\begin{tabularx}{\textwidth}{p{39mm}X}
\toprule
符号 & 含义 \\
\midrule
$\mathbf{x}$ & 决策变量或状态向量；各分量单位在首次定义处说明 \\
$\theta$ & 模型参数向量；填写参数含义、单位与依据 \\
$\mathcal{D}$ & 本问实际使用的数据集合及其支持范围 \\
$\mathcal{X}$ & 满足本问约束的可行域 \\
$f$ & 与本问目标对应的函数；量纲由实际定义确定 \\
\bottomrule
\end{tabularx}
\end{table}
\FloatBarrier
\renewcommand{\thetable}{1-\arabic{table}}
\setcounter{table}{0}
\templatehint{这些符号只演示符号表版式，不能据此预设后续题目的变量或模型。}

\section{模型建立与求解}
% 本处集中管理题数及顺序；增加一问时复制现有q文件并增加一条input。
\startquestion{1}{一}{填写本问研究对象与应直接回答的问题。}
\subsubsection*{1）模型的准备}
先界定研究对象、输入数据及单位，说明与前问或共用资料的接口。介绍关键变量和参数如何取得，再解释需要解决的困难。

\noindent\hspace*{2em}\topic{资料与变量}
说明实际可用信息、处理方式和时间空间支持。无法识别的量保留未知，不用文字补出尚未取得的数据。

\subsubsection*{2）模型的建立}
从本问目标引出数学关系。式~\eqref{eq:format-demo}仅演示公式排版、编号与引用，正式稿应替换成实际模型或删除。
\begin{equation}
\min_{\mathbf{x}\in\mathcal{X}} f(\mathbf{x};\theta).
\label{eq:format-demo}
\end{equation}
式中，$\mathbf{x}$表示变量，$\theta$表示参数，$\mathcal{X}$表示可行域，$f$表示目标函数。此处只说明数学记号，不意味着本题应使用优化模型。

\subsubsection*{3）模型的求解与分析}
说明实际采用的求解步骤、停止条件和必要核查；结果先回答本问，再解释关键关系。表~\ref{tab:q1-demo}演示三线表，不含研究数值。
\begin{table}[htbp]
\centering
\caption{本问结果展示位置（格式示例）}\label{tab:q1-demo}
\fontsize{9.5}{12}\selectfont
\setlength{\tabcolsep}{4pt}
\begin{tabularx}{\textwidth}{p{25mm}X X}
\toprule
项目 & 正式结果 & 对应解释 \\
\midrule
主要答案 & 填写已确认结果 & 说明它如何回答本问 \\
必要核查 & 填写实际核查结果 & 说明支持范围与未解项 \\
\bottomrule
\end{tabularx}
\end{table}
\FloatBarrier

\startquestion{2}{二}{填写本问的研究目标，以及相对前问新增的条件。}
\subsubsection*{1）模型的准备}
说明本问实际继承了前问哪些有效资料、定义和结果，以及它们的单位、支持范围与必要限制。新增资料应有明确来源，不能只因为题号相邻就假定所有输出直接可用。

\subsubsection*{2）模型的建立}
围绕本问目标，说明模型选择理由、关键数学关系、参数依据和适用条件。由本问研究审核与你共同确定具体方法；此处没有预设算法，也没有替本问完成研究。

\subsubsection*{3）模型的求解与分析}
填写已确认的方法如何实现和核查，再按本问要求组织真实表图与解释。关键数值从正式结果取用；条件、区间、概率或未达状态应有清晰定义。

本问结束时概括已经回答什么、仍缺什么，并说明其对下一问的数据、假设和整体方向有什么影响。实际不需要该段时按论述自然合并，不机械保留模板提示。
\FloatBarrier

\startquestion{3}{三}{填写本问的研究目标，以及相对前问新增的条件。}
\subsubsection*{1）模型的准备}
说明本问实际继承了前问哪些有效资料、定义和结果，以及它们的单位、支持范围与必要限制。新增资料应有明确来源，不能只因为题号相邻就假定所有输出直接可用。

\subsubsection*{2）模型的建立}
围绕本问目标，说明模型选择理由、关键数学关系、参数依据和适用条件。由本问研究审核与你共同确定具体方法；此处没有预设算法，也没有替本问完成研究。

\subsubsection*{3）模型的求解与分析}
填写已确认的方法如何实现和核查，再按本问要求组织真实表图与解释。关键数值从正式结果取用；条件、区间、概率或未达状态应有清晰定义。

本问结束时概括已经回答什么、仍缺什么，并说明其对下一问的数据、假设和整体方向有什么影响。实际不需要该段时按论述自然合并，不机械保留模板提示。
\FloatBarrier

\startquestion{4}{四}{填写本问的研究目标，以及相对前问新增的条件。}
\subsubsection*{1）模型的准备}
说明本问实际继承了前问哪些有效资料、定义和结果，以及它们的单位、支持范围与必要限制。新增资料应有明确来源，不能只因为题号相邻就假定所有输出直接可用。

\subsubsection*{2）模型的建立}
围绕本问目标，说明模型选择理由、关键数学关系、参数依据和适用条件。由本问研究审核与你共同确定具体方法；此处没有预设算法，也没有替本问完成研究。

\subsubsection*{3）模型的求解与分析}
填写已确认的方法如何实现和核查，再按本问要求组织真实表图与解释。关键数值从正式结果取用；条件、区间、概率或未达状态应有清晰定义。

本问结束时概括已经回答什么、仍缺什么，并说明其对下一问的数据、假设和整体方向有什么影响。实际不需要该段时按论述自然合并，不机械保留模板提示。
\FloatBarrier


\FloatBarrier
\section{模型评价、改进与推广}
\subsection{模型的优点}
围绕模型的具体性质和已有证据说明优点，例如解释性、约束处理或计算效率。不要只罗列“准确”“稳定”“创新”等缺少依据的评价。

\subsection{模型的缺点}
集中说明实际存在的资料、识别、假设或验证限制，并指出它们影响哪些结论。必要的条件和不利事实仍在相关结果处准确说明。

\subsection{模型的改进与推广}
提出与本次实际缺口对应的改进路径，以及更换资料、参数或场景后需要重新核查的内容。不把尚未实施的改进写成已有成绩。
\FloatBarrier

\stepcounter{section}
\renewcommand{\refname}{\texorpdfstring{\chinese{section}、\quad 参考文献}{参考文献}}
\begin{thebibliography}{99}
\small
\raggedright
\bibitem{template-guide} 本模板使用说明[Z]. 仅演示文献编号与引用位置，不作为正式研究来源；成稿时替换并核实原始出处。
\end{thebibliography}

\ifshowdisclosure\noindent\textbf{工具使用声明（格式位置示例）：}按当届比赛规定填写本次真实使用的工具、用途、日期和人工核查情况。本模板不提供默认工具名称、模型名称或既往使用历史。
\fi
\ifshowappendix\FloatBarrier
\clearpage
\section{附录}
\renewcommand{\thefigure}{A\arabic{figure}}\setcounter{figure}{0}
\renewcommand{\thetable}{A\arabic{table}}\setcounter{table}{0}
\subsection*{附录1：补充推导与详细结果}
按正文需要放置补充推导、详细表格或其他可检查证据。附录并非默认存放全部数据；范围由题目、官方要求与本次交付决定。

\subsection*{附录2：程序说明与源码}
说明运行入口、输入输出与必要依赖。若要求全部正式源码，按确切版本完整列出；若不要求全量，依据本次约定组织内容。
\begin{Verbatim}[fontsize=\small,breaklines=true,breakanywhere=true]
# 这里替换为本次比赛正式程序或输入源码的命令。
# 这段文字只是代码附录的格式位置示例，不是可运行分析。
\end{Verbatim}
\templatehint{本样张保留文后声明与附录位置。config.tex中的开关可控制是否显示，实际要求以当届规则为准。}
\fi
\end{document}
